\documentclass[11pt]{article}
\usepackage[a4paper,margin=25mm]{geometry}
\usepackage{amsmath,amssymb,amsthm}
\usepackage[hidelinks]{hyperref}
\newtheorem{theorem}{Theorem}
\newtheorem{proposition}[theorem]{Proposition}
\theoremstyle{remark}
\newtheorem*{remark}{Remark}
\title{Subdividing an Edge of $K_9$:\\A Disproof of Conjecture 00000003476}
\author{GitHub: gaochengzhecpu}
\date{}
\begin{document}
\maketitle
\begin{abstract}
The source asserts that subdividing an edge decreases the Wiener index
if and only if the betweenness of that edge is below a quarter of the
number of vertices. In the complete graph $K_9$ every edge has
betweenness $1<9/4$, yet subdividing an edge raises the Wiener index
from $36$ to $53$. The graph, the subdivision, the Wiener index, the
betweenness and the negation of the criterion are formalised in Lean.
In fact subdivision never decreases the Wiener index of a connected
graph. An earlier submission argued with $K_5$ but had none of these
objects in Lean; its error is described below.
\end{abstract}

\paragraph{Definitions and reading of the source.}
Let $G$ be a finite connected simple graph with $n$ vertices and
distance function $d$. Its Wiener index is
$W(G)=\sum_{\{x,y\}}d(x,y)$, the sum over unordered pairs of vertices.
Subdividing an edge $e=ab$ gives the graph $G_e$: the edge $ab$ is
removed and a new vertex $w$ adjacent to exactly $a$ and $b$ is added.
The betweenness of $e$ is
\[
 B(e)=\sum_{\{s,t\}}\frac{\sigma_{st}(e)}{\sigma_{st}},
\]
where $\sigma_{st}$ is the number of shortest paths from $s$ to $t$ and
$\sigma_{st}(e)$ the number of those through $e$. The source claims, for
every such $G$ and $e$,
\begin{equation}\label{eq:crit}
 W(G_e)<W(G)\iff B(e)<\frac n4 .
\end{equation}

\begin{theorem}
Let $G=K_9$ and let $e$ be an edge. Then $B(e)=1<9/4$, while
$W(G)=36$ and $W(G_e)=53$. Hence \eqref{eq:crit} fails.
\end{theorem}
\begin{proof}
Any two distinct vertices $s,t$ of $K_9$ are adjacent, so the only
shortest path between them is the edge $st$. It passes through $e$
exactly when $\{s,t\}=e$. Thus $B(e)=1$. All $\binom92=36$ distances
equal $1$, so $W(G)=36$.

Let $e=ab$. The graph $G_e$ has $10$ vertices and is connected, so each
of its $\binom{10}2=45$ distances is at least $1$. Already this gives
$W(G_e)\ge45>36$. Precisely, the pairs at distance $2$ are $\{a,b\}$ and
$\{w,x\}$ for the seven vertices $x\ne a,b$; all other pairs are
adjacent. So $W(G_e)=45+8=53$.
\end{proof}

\begin{proposition}
For every finite connected graph $G$ and every edge $e$,
$W(G_e)>W(G)$.
\end{proposition}
\begin{proof}
Let $e=ab$ and let $d'$ be the distance in $G_e$. A walk in $G_e$
between two old vertices $x,y$ yields a walk in $G$ that is no longer:
replace each passage $a,w,b$ or $b,w,a$ by the edge $ab$, and delete
each passage $a,w,a$ or $b,w,b$. Hence $d'(x,y)\ge d(x,y)$. The pairs
containing $w$ contribute a positive amount. Summing gives
$W(G_e)>W(G)$.
\end{proof}

So the left side of \eqref{eq:crit} never holds, and the criterion fails
for every edge with $B(e)<n/4$.

\begin{remark}
The value of $B(e)$ depends on the convention only mildly. Summing over
ordered pairs gives $2$; omitting the pair of endpoints gives $0$;
normalised versions are smaller. All are below $9/4$, which is why
$K_9$ is used and not a smaller complete graph. If the quarter is taken
of the $10$ vertices of $G_e$, the same inequalities hold. For the
clause on bipartite graphs, $K_{7,7}$ has $n=14$ and
$B(e)=1+2\cdot\tfrac67=\tfrac{19}7<\tfrac{14}4$ for every edge, so by
the proposition the criterion fails there as well. The proposition and
this last computation are not formalised.
\end{remark}

\section*{The earlier submission and its error}
Pull request 225 of the competition repository (by
\texttt{orionsheep}) proposed the same kind of counterexample with
$K_5$. It was closed without merging on 3 October 2026. The reviewer's
stated reason was that the main Lean theorem was arithmetic on a
hard-coded list of distances, and that no graph, Wiener index,
subdivision or betweenness was defined in Lean.

That is the error: the Lean file checked that a list of numbers sums to
$19$ and that $19>10$, which says nothing about graphs. There is also a
gap in the example itself. With betweenness summed over ordered pairs
an edge of $K_5$ has betweenness $2>5/4$, so $K_5$ refutes the
criterion only under some conventions.

The present submission defines the four objects in Lean for arbitrary
finite graphs, states the criterion with a universal quantifier, and
uses $K_9$, for which the ordered convention is covered too.

\section*{Formal verification and scope}
Graphs are Mathlib's \texttt{SimpleGraph} and distances are Mathlib's
\texttt{SimpleGraph.dist}. \texttt{subdivide G a b} is a graph on
\texttt{Option V}; the new vertex is \texttt{none}.
\texttt{wiener} is half the sum of distances over ordered pairs.
\texttt{sigma} and \texttt{sigmaThrough} count the walks whose length
equals the distance, and those among them whose edge list contains
$e$; \texttt{betweenness} is the resulting sum over unordered pairs and
\texttt{betweennessOrdered} the sum over ordered pairs.

For a complete graph on any vertex type, \texttt{sigma\_top} and
\texttt{sigmaThrough\_top} prove $\sigma_{st}=1$ and
$\sigma_{st}(e)\in\{0,1\}$ according to whether $\{s,t\}=e$, by
classifying all walks of length one. This gives
\texttt{betweenness\_K9}: $B(e)=1$, and
\texttt{betweennessOrdered\_K9}: the ordered sum is $2$.
\texttt{wienerTwice\_K9} computes the ordered distance sum $72$.
\texttt{subdivide\_K9\_connected} proves that the subdivided graph is
connected, and \texttt{wienerTwice\_ge} proves that a connected graph
on $m$ vertices has ordered distance sum at least $m(m-1)$. Hence
\texttt{wiener\_increases}: $W(K_9)<W((K_9)_e)$. Lean proves the bound
$W((K_9)_e)\ge45$, which is what the disproof needs; the exact value
$53$ is the hand computation above.

The proposition \texttt{ClaimedCriterion} is \eqref{eq:crit} quantified
over all finite connected graphs and edges. Its negation is the theorem
\begin{center}\small\texttt{conjecture\_false}.\end{center}
The theorem
\begin{center}\small\texttt{conjecture\_false\_ordered}\end{center}
is the same for the ordered convention. No \texttt{sorry},
\texttt{native\_decide} or custom axiom occurs.

The disproof concerns the stated equivalence. No claim is made about
the phrase on the cut-load inequality chain, which is not a precise
statement.

\section*{Source and reproduction}
The exact bilingual source is included byte-for-byte as
\texttt{SOURCE.md}. Its repository path is
\begin{center}\small\texttt{conjectures/00000003476.md}.\end{center}
The project pins Lean 4.19.0 and Mathlib commit
\begin{center}\small\texttt{c44e0c8ee63ca166450922a373c7409c5d26b00b}.\end{center}
From \texttt{lean/}, run \texttt{lake build}, then
\begin{center}\small\texttt{lake env lean -DwarningAsError=true Main.lean}.\end{center}
The included public-Git manifest locks all dependency revisions. The
\texttt{verification/} directory records the fresh-build logs,
theorem-axiom reports, artifact hashes, review and final PDF
inspection. No supplementary numerical computation is required.
\end{document}
